\section{Retrieval does not imply appropriate use}\label{sec:C14:use}
Memory can fail through excessive influence as well as omission. MemSyco-Bench constructs decisions where relevant historical material must be rejected as factual evidence, bounded to its original scope, reconciled with current evidence, updated, or used for legitimate personalization~\cite{cumb2026memsycov3}. These categories distinguish retrieval success from epistemically appropriate use. A stored preference is evidence about a user's preference; it does not become evidence that the preferred claim is objectively correct. A discriminating counterfactual preserves the current evidence and changes only the historical preference. An objective answer should remain invariant, while a preference-dependent recommendation may appropriately change. That conditional distinction is stronger than rewarding either universal obedience to memory or universal suppression.

Response concentration is a separate outcome. The October MemBubble benchmark derives preferences from real conversations and constructs topically related queries in distinct scenarios, then measures diversity using functional-equivalence classes of sampled answers~\cite{tan2026membubble}. Its query filtering and diversity pre-test define a selected evaluation population; the result cannot estimate all deployment queries without a sampling argument. An original evaluation recommendation is to report valid-personalization accuracy and response-distribution concentration jointly. Higher entropy can include irrelevant alternatives, while lower entropy can reflect a justified constraint. Hold the decoder and sampling budget fixed, inspect the equivalence-class judge, and separate memory-induced narrowing from changes already implied by the current request.

An agent can store a fact, retrieve it when asked directly, and fail to apply it when relevance is implicit. InMind isolates this distinction with paired storage and in-context controls and an answer-blind target-recall measurement~\cite{inmind2026}. Its 125-task diagnostic reports a large gap between direct recall and application under indirect queries. The result concerns the tested systems and selected implicit associations, not every possible retrieval mechanism. Missing broad negative controls also matters when evaluating whether a remedy overapplies memory.

The mechanism is a sequencing problem. A query-conditioned retriever must decide what is relevant before the full policy sees the candidate fact. If relevance depends on a bridge supplied by world knowledge, surface similarity can miss it. Increasing embedding dimensionality does not necessarily solve the need to infer that bridge. A structured profile, proactive constraint store, query expansion, or model-assisted memory search can change the access interface, but each introduces a different cost and failure mode.

The official InMind artifact makes the evaluation boundary inspectable. Its timeline builder constructs the ordered conversation for a selected task, while its judge-payload builder separates measurements through metric-specific payloads~\cite{repo_inmind2026}. These are benchmark construction and scoring paths rather than a released solution to the memory-use problem. A reproduction must preserve the paired controls, target-fact placement, answer-blind measurement, and system output schema. Omitting a control can turn a retrieval diagnosis into an ambiguous end-to-end failure.

InMind's reported in-context control reaches 84.0\%, while the six evaluated memory systems reach at most 14.4\% on indirect application and can reach 100\% on direct recall~\cite{inmind2026}. The separately evaluated dense NaiveRAG baseline reaches 16.0\% on indirect application; it is not included in that six-system maximum. These percentages have different experimental conditions and should not be subtracted as if they were a randomized treatment effect. The official timeline code also freezes the post-history state for direct and indirect queries, preventing one query's answer from teaching the subsequent query the target fact. This control is crucial: sequentially asking the direct question first would change the state and confound the diagnosis.

Appropriate use also requires authority calibration after retrieval. A memory can be relevant as background evidence without being a binding constraint. A tentative preference should not override a current explicit instruction, and an old successful procedure should not override a changed tool contract. One proposed interface labels each item with its role: observation, preference, candidate procedure, or current obligation. The model still interprets the item, while deterministic controls can prevent some role promotions at consequential boundaries.

Formally, let $z$ denote the current task state and $m$ a retrieved item. The useful question is the change in expected loss caused by making $m$ visible, conditional on its status and provenance. Similarity is a retrieval signal; causal usefulness is an outcome under an intervention. Estimating the latter requires matched continuations or randomized memory exposure where appropriate. Observational correlation between retrieval and success is confounded because easier or familiar tasks can both retrieve better and succeed more often.

An evaluation should include four conditions: the fact absent, present in storage, explicitly retrieved, and permanently visible as a diagnostic control. Compare action correctness and inappropriate-use rates, not only answer recall. If explicit exposure repairs performance, the bottleneck is likely upstream of use under those conditions. If it does not, missing bridge knowledge or interpretation remains possible. This decomposition gives a concrete experimental route from an aggregate memory score to a mechanism-specific conclusion.

\subsection{Identifying memory influence}
A controlled exposure study can make the causal estimand explicit. Let $X_i$ bind a frozen task, memory-bank snapshot and model configuration. Let $B_i\in\{0,1\}$ indicate whether one eligible memory item is exposed, with known assignment probability $e_i=P(B_i=1\mid X_i)\in(0,1)$. Write $L_i(1)$ and $L_i(0)$ for the two potential decision losses and observe $L_i=B_iL_i(1)+(1-B_i)L_i(0)$. Under randomized assignment conditional on $X_i$, consistent exposure, and no cross-trial state interference,
\begin{equation}
 \widehat\Delta_n=\frac1n\sum_{i=1}^n
 \left[\frac{B_iL_i}{e_i}-\frac{(1-B_i)L_i}{1-e_i}\right]
\end{equation}
has expectation $n^{-1}\sum_i\mathbb E[L_i(1)-L_i(0)]$, the average exposure effect for the sampled task population. The identity follows by conditioning each term on $X_i$ and its potential losses. Negative $\Delta$ means beneficial exposure on this evaluated distribution. Very small propensities can make the estimator unstable, so uncertainty and the assignment range must be reported.

This construction is a proposed experiment, not a measured memory benefit. Use sandboxed or read-only continuations with reset state where withholding evidence is permissible; current authorization and mandatory obligations are not randomized away. A shared persistent bank updated by earlier trials breaks the no-interference premise unless the trajectory itself is the randomized unit. Delayed or missing labels require an additional observation model: logging admission propensity alone does not identify outcomes for permanently censored facts. The experiment also distinguishes the influence of an item from its factual correctness and from the quality of the writer that produced it.
